\section{Quivers}
In this section, we introduce the point count polynomial $R(Q;q)$ of a \Louise %
 quiver, and study its basic properties.  We use this polynomial to define the \emph{$Q$-Catalan number} of $Q$, an integer that we conjecture to be nonnegative in \cref{conj:Catalan}. We refer the reader to~\cite{FWZ_book} for an excellent introduction to quivers and cluster algebras.

\subsection{Quiver mutation}\label{sec:quiver}
Recall that a \emph{quiver} is a directed graph without directed cycles of length $1$ and $2$.

Given a quiver $Q$ and a vertex $j\in V(Q)$, one can define another quiver $Q'=\mu_j(Q)$ called the \emph{mutation of $Q$ at $j$}. This operation preserves the set of vertices:  $V(Q')=V(Q)$, and changes the set of arrows as follows:
\begin{itemize}
\item for every length $2$ directed path $i\to j\to k$ in $Q$, add an arrow $i\to k$ to $Q'$;
\item reverse all arrows in $Q$ incident to $j$;
\item remove pairs of opposite arrows (i.e.  directed $2$-cycles) in the resulting directed graph, one pair at a time, until no such pairs are present.
\end{itemize}

An \emph{ice quiver} is a quiver $\Qice$ whose vertex set $\Vice=V(\Qice)$ is partitioned into \emph{frozen} and \emph{mutable} vertices: $\Vice=\Vfro\sqcup \Vmut$. We automatically omit all arrows in $\Qice$ both of whose endpoints are frozen. For an ice quiver $\Qice$, its \emph{mutable part} is the induced subquiver of $\Qice$ with vertex set $\Vmut$.  %
 For a set $S\subset \Vmut$, let $\Qicefr[S]$ denote the ice quiver obtained from $\Qice$ by further declaring all vertices in $S$ to be frozen. We write $\Qice-S$ for the ice quiver obtained from $\Qice$ by removing the vertices in $S$.  For a simply-laced Dynkin diagram $D$, we say that $Q$ or $\Qice$ has type $D$, or is a $D$-quiver, if the underlying graph of $Q$ is isomorphic to $D$.


 Let $\Qice$ be an ice quiver with $\Vice=[n+m]$, $\Vmut=[n]$, and $\Vfro=[n+1,n+m]:=\{n+1,n+2,\dots,n+m\}$. We represent it by an $(n+m)\times n$ \emph{exchange matrix} $\BQice=(b_{ij})$, defined by
 \begin{equation*}%
   b_{ij}=\#\{\text{arrows $i\to j$ in $\Qice$}\} - \#\{\text{arrows $j\to i$ in $\Qice$}\}.
 \end{equation*}
The top $n\times n$ submatrix $\BQ$ of $\BQice$ is skew-symmetric, and is called the \emph{principal part} of $\BQice$.  
%

%
%
%
%
%
%
%
%
%
%
%
%
%
%
%
%
 
%

Let $\Bice$ be an $(n+m)\times n$ matrix. We denote $\cork(\Bice):=n-\rk(\Bice)$. Following~\cite{LS}, we say that an $(n+m)\times n$ matrix $\Bice$ is \emph{really full rank} if the rows of $\Bice$ span $\Z^n$ over $\Z$. We say that $\Qice$ is \emph{really full rank} if $\BQice$ is really full rank, and write $\cork(\Qice):=\cork(\BQice)$.   We say that $\Qice$ is \emph{torsion-free} if the abelian group $\Z^{n+m}/\BQice \Z^n$ is torsion-free, equivalently, if $\Z^n/\BQice^T \Z^{n+m}$ is torsion-free.  Thus, $\Qice$ is really full rank if and only if it is both full rank and torsion-free.

\begin{definition}
Given a quiver $Q$, we say that an ice quiver $\Qice$ with mutable part $Q$ is a \emph{minimal extension} of $Q$ if $\Qice$ has $m=\cork(Q)$ frozen vertices and is really full rank.
\end{definition}

\begin{lemma}\label{lem:ext}
If $Q$ is torsion-free then it admits a minimal extension.
%
\end{lemma}
\begin{proof}
If $Q$ is torsion-free, we can find vectors $d_1,\ldots,d_m \in \Z^n$ that form a basis of $\Z^n/B(Q)^T \Z^n$.  Define $\Qice$ with $m$ frozen vertices so that the bottom $m$ rows of the exchange matrix $\Bice$ are equal to $d_1,\ldots,d_m$.
\end{proof}
%
%

\begin{example}
The exchange matrices of the two quivers $Q,Q'$ in \figref{fig:LA}(c) are given by
$$
\BQ=\begin{pmatrix}
0 & 2 & -2 \\
-2&0&2\\
2&-2&0
\end{pmatrix} \qquad \text{and} \qquad
B(Q')=\begin{pmatrix}
0 & -1& -1 &2 \\
1&0&1 &-1\\
1&-1&0 &-1 \\
-2&1&1&0
\end{pmatrix}.
$$
Thus, $\cork(Q)=1$, $\cork(Q')=0$, and neither quiver is torsion-free.
\end{example}

We say that an ice quiver $\Qice$ is \emph{isolated} if it has no arrows. For example, the quiver in \figref{fig:simple}(c) is isolated.

We call an ice quiver $\Qice$ \emph{acyclic} if it has no directed cycles, and \emph{mutation acyclic} if it is mutation equivalent to an acyclic ice quiver.  For example, any quiver of Dynkin type is acyclic, as is the quiver in \figref{fig:LA}(a).  The quiver in \figref{fig:simple}(a) and the left quiver in \figref{fig:LA}(b) are not acyclic, but they are both mutation acyclic, since mutating a single vertex in each of these quivers makes them acyclic.

An edge $u\to v$ in a quiver $Q$ is called a \emph{separating edge}~\cite{Muller} if it does not belong to a bi-infinite walk in $Q$. Here, a \emph{bi-infinite walk} is a sequence $(w_j)_{j\in\Z}$ of vertices in $Q$ such that for each $j\in\Z$, $Q$ contains an arrow $w_j\to w_{j+1}$.

We define the class of \emph{\Louise quivers}, called ``Louise" in \cite{MSLA,LS}, as follows.
\begin{itemize}
\item Any isolated quiver $Q$ is \Louise.
\item Any quiver that is mutation equivalent to a \Louise quiver is \Louise.
\item Suppose that a quiver $Q$ has a separating edge $u\to v$, and that all three quivers $Q-\{u\}$, $Q-\{v\}$, $Q-\{u,v\}$ are \Louise. Then $Q$ is \Louise.
\end{itemize}
We say that an ice quiver $\Qice$ is \emph{\Louise} if its mutable part $Q$ is \Louise.    See \cref{fig:LA}.
\begin{figure}
  \includegraphics[width=\textwidth]{figures/Louise}
  \caption{\label{fig:LA} Acyclic, locally acyclic and non-\Louise quivers.}
\end{figure}

\subsection{Cluster algebras}
Let $\Qice$ be an ice quiver with $\Vice=[n+m]$, $\Vmut=[n]$, and $\Vfro=[n+1,n+m]$.  We associate to $\Qice$ the initial seed $(x_1,x_2,\ldots,x_{n+m})$ of \emph{cluster variables}, considered as elements in the function field $\Fcal = \Q(x_1,x_2,\dots,x_{n+m})$.  The variables $x_{n+1},\ldots,x_{n+m}$ are called \emph{frozen variables}.  For a mutation $Q' = \mu_j(Q)$, we associate the seed $(x'_1,x'_2,\ldots,x'_{n+m})$, where $x'_i = x_i$ if $i \neq j$ and one has the new cluster variable
$$
x'_j = \frac{\prod_{r \to j} x_r + \prod_{j \to s} x_s}{x_j} \in \Fcal.
$$
By repeatedly mutating, we generate (possibly infinitely) many seeds and cluster variables. 
 We denote by $\AQice=\ABQice$ the cluster algebra associated to $\Qice$. This is the $\Z$-subalgebra of $\Q(x_1,x_2,\dots,x_{n+m})$ generated by all cluster variables and the inverses of frozen variables. 
 
The \emph{cluster variety} $\XQice=\XBQice$ is defined to be the scheme
\begin{equation*}%
\XQice:=\Spec(\AQice).
\end{equation*}
We say that $\AQice$ (resp. $\XQice$) is a cluster algebra (resp. cluster variety) \emph{of type $Q$}, where $Q$ is the mutable part of $\Qice$.  We say that $\AQice$ is isolated, acyclic, (really) full rank, if $\Qice$ is isolated, mutation acyclic, (really) full rank, respectively.  If $\Qice$ is a \Louise quiver, then $\AQice$ is locally acyclic in the sense of Muller~\cite{Muller}.


\begin{proposition}[{\cite[Theorem 7.7]{Muller}, \cite[Theorem 10.1]{LS}}]\label{prop:Mulsmooth}
Suppose that $\Qice$ is \Louise and really full rank.  Then $\XQice_\C$ is a smooth complex algebraic variety and $\XQice_{\bar \F_q}$ is smooth for any prime power $q$.
\end{proposition}

\begin{proposition}[{\cite[Corollary 5.4]{Muller}}]\label{prop:Mulcover}
Let $i \to j$ be a separating edge in $\Qice$.  Then the open sets $\{x_i \neq 0\}$ and $\{x_j \neq 0\}$ cover $\XQice$.
\end{proposition}

\begin{proposition}[{\cite[Proposition 3.1, Lemma 3.4, Theorem 4.1]{Muller}}]\label{prop:Mulloc}
Suppose that $i \in V(\Qice)$ is a mutable vertex such that $Q- \{i\}$ is \Louise.  Then $\AQice[x_i^{-1}] \simeq \AX(\Qice[i])$.
\end{proposition}

\subsection{Quiver point count}
For a mutable quiver $Q$, we define the function $\RQ(q)$ as follows. Choose a really full rank ice quiver $\Qice$ with mutable part $Q$ and $m$ frozen vertices. Then for $q$ a prime power, we set
\begin{equation}\label{eq:RQ_dfn}
  \RQ(q):=\frac{\#\XQice(\F_q)}{(q-1)^m}.
\end{equation}
For a rational function $R(q)=P(q)/Q(q)$, we let the \emph{degree} $\deg(R)$ be the difference $\deg(P)-\deg(Q)$. 
\begin{proposition}[{\cite[Proposition 5.11 and Theorem 10.5]{LS}}]\label{prop:Q_pcnt}
  Let $Q$ be a quiver with $n$ vertices.
  \begin{enumerate}
  \item The function $\RQ(q)$ defined in~\eqref{eq:RQ_dfn} does not depend on the choice of $\Qice$.
  \item Suppose $Q$ is \Louise. Then $\RQ(q)$ is a rational function in $q$ of degree $n$. %
  \end{enumerate}
\end{proposition}
If $u \to v$ is a separating edge in $Q$ and $Q-\{u\},Q-\{v\},Q-\{u,v\}$ are \Louise, then we have the recurrence
\begin{equation}\label{eq:RQ_recurrence}
  \RQ(q) = (q-1)\RQX{Q-\{u\}}{q}+ (q-1)\RQX{Q-\{v\}}{q}- (q-1)^2\RQX{Q-\{u,v\}}{q},
\end{equation}
which follows from Propositions~\ref{prop:Mulcover} and~\ref{prop:Mulloc}.

\begin{proposition}[{\cite[Proposition 3.9]{LS2}}] \label{prop:acycliccount}
Let $Q$ be an acyclic quiver with $n$ vertices.  Then 
$$
\RQ(q) = \sum_{k \geq 0} a_k (q-1)^{n-2k} q^k,
$$
where $a_k$ is the number of independent sets of size $k$ in the underlying undirected graph of~$Q$.
\end{proposition}

When $Q$ itself is really full rank, $\RQ(q)$ is a polynomial in $q$. However, $\RQ(q)$ is in general a genuine rational function whose denominator is a power of $(q-1)$. For example, for $Q$ a single isolated vertex (we denote this quiver by $A_1$; it has corank $1$), we have
\begin{equation}\label{eq:A1}
  \RQX{A_1}{q}=\frac{q^2-q+1}{q-1}.
\end{equation}
For an ice quiver $\Qice$ with mutable part $Q$ and $m$ frozen vertices, we set
\begin{equation*}%
  \RQice(q):=(q-1)^m\RQ(q).
\end{equation*}
Thus, when $\Qice$ is really full rank and \Louise, $\RQice(q)$ is a polynomial in $q$ of degree $n+m$ satisfying $\#\XQice(\F_q)=\RQice(q)$ for all prime powers $q$.
\begin{definition}\label{def:QCat}
Let $Q$ be a torsion-free quiver, and $\Qice$ be a minimal extension as in \cref{lem:ext}.  If $\RQice(q)$ is a polynomial in $q$, define %
  \begin{equation*}%
    \chiQ:=\RQice(1).
  \end{equation*}
  By \cref{prop:Q_pcnt}, $\chiQ$ does not depend on the choice of $\Qice$.
\end{definition}
We call $\chiQ$ the \emph{Catalan number} of the quiver $Q$, or simply the \emph{$Q$-Catalan number}.

\begin{conjecture}\label{conj:Catalan}
Suppose that $Q$ is a \Louise quiver. %
Then $\chiQ \geq0$. %
\end{conjecture}
%

We note that $\chiQ$ can be zero.  For example, let $Q$ be an acyclic orientation of the three-cycle, and let $\Qice$ be a minimal extension. Thus, $\Qice$ is a really full rank quiver with $\cork(Q)=1$ frozen vertices.
%
 Then by \cref{prop:acycliccount}, we have $R(Q;q)=(q-1)^3+3(q-1)q$, and thus $R(\Qice;q) = (q-1)^4+3(q-1)^2q$ and  $\chiQ =0$.  For another example, the quiver $Q$ in \figref{fig:LA}(a) satisfies $\cork(Q) = 0$ and $\chiQ = 0$.  %

Consider a polynomial $R(q)$ of degree $n$ with integer coefficients  that is palindromic (i.e.  satisfies $q^nR(1/q)=R(q)$). One can show by induction that it can be written as $\sum_{k=0}^{\lfloor n/2\rfloor} a_k (q-1)^{n-2k} q^k$ for unique coefficients $a_0,a_1,\dots,a_{\lfloor n/2\rfloor}\in\Z$. It follows from the \emph{curious Lefschetz property} shown in~\cite{LS} that the point count polynomial of any \Louise really full rank quiver is palindromic. We thank the anonymous referee for suggesting the first part of the following strengthening of \cref{conj:Catalan}. 
\begin{conjecture}
Let $Q$ be a \Louise quiver, and let $\Qice$ be a really full rank quiver with $n$ vertices and mutable part $Q$. Then we have 
\begin{equation*}%\label{eq:*}
  \RQice(q) = \sum_{k=0}^{\lfloor n/2\rfloor} a_k (q-1)^{n-2k} q^k, \quad\text{where each $a_k \geq0$.}
\end{equation*}
Moreover, the cluster variety $\XQice$ over a field $\F$ may be represented as a disjoint union over $0\leq k\leq\lfloor n/2\rfloor$ of $a_k$ copies of $(\F^\ast)^{n-2k} \times\F^k$.
\end{conjecture}
